% ===================== ANEXOS =====================
\appendix

\section{Registro HAZOP + What-If}
\label{appendix:hazop}

Registro completo de sesiones HAZOP/What-If con nodos, palabras guía, desviaciones, causas, consecuencias, salvaguardas existentes, calificación de riesgo y recomendaciones. Fuente: archivo \texttt{data/hazop.csv}.

\section{Memoria de cálculo hidráulico}
\label{appendix:hidraulica}

Memoria completa del análisis nodal NFPA 13 del área hidráulicamente más remota, incluyendo tubería por tubería, rociador por rociador y la curva de suministro vs.\ demanda. Los resultados se generan con \texttt{src/pci/hidraulica.py} y se reproducen en \texttt{notebooks/01\_hidraulica.ipynb}.

\section{Memoria de cálculo — agentes de extinción}
\label{appendix:agentes}

Cálculos de masa de agente limpio (NFPA 2001), tasa de aplicación de espuma (NFPA 11/16) y/o masa de CO\textsubscript{2} (NFPA 12) según el sistema seleccionado. Módulos: \texttt{src/pci/agentes\_limpios.py}, \texttt{src/pci/espumas.py}.

\section{Detección y alarma — planillas NFPA 72}
\label{appendix:deteccion}

Planillas de dispositivos, cálculos de batería, caída de tensión en NAC/SLC y niveles de audibilidad. Módulos: \texttt{src/pci/deteccion.py}, \texttt{src/pci/alarma\_nac.py}.

\section{Clasificación de áreas NEC}
\label{appendix:clasificacion}

Listado de áreas clasificadas con Clase/División (o Zona), grupo de material, código de temperatura y fuente de ignición. Fuente: \texttt{data/clasificacion\_areas.csv}.

\section{Planos e isométricos}
\label{appendix:planos}

\begin{itemize}
\item P\&ID general del sistema PCI (símbolos ISA 5.1 + NFPA 170).
\item Isométricos hidráulicos por rama.
\item Planta de distribución de detectores y notificadores.
\item Diagramas unifilares eléctricos y planos de áreas clasificadas NEC.
\end{itemize}

\section{Listado de equipos}
\label{appendix:equipos}

\begin{table}[htbp]
\centering
\caption{Listado maestro de equipos PCI}
\label{tab:equipos}
\setcounter{itemcount}{0}
\begin{tabularx}{\textwidth}{@{}NlXcl@{}}
\toprule
\multicolumn{1}{c}{\textbf{Item}} & \textbf{Tag} & \textbf{Descripción} & \textbf{Cant.} & \textbf{Norma} \\
\midrule
& FP-001 & Bomba contra incendio principal & 1 & NFPA 20 \\
& FP-002 & Bomba jockey & 1 & NFPA 20 \\
& TK-001 & Tanque de reserva & 1 & NFPA 22 \\
& FACP-01 & Panel de alarma central & 1 & NFPA 72 \\
\bottomrule
\end{tabularx}
\end{table}
